\chapter{总结与展望}
\label{cha:conclusion}

\section{工作总结}
\label{sec:summary}

本文提出了一种基于 Attention-MAPPO 的多无人机 MEC 轨迹与任务卸载协同优化框架。通过将轨迹控制和请求级卸载解耦为两个协同决策层，该框架降低了直接端到端联合动作建模的结构复杂度，并提高了决策过程的可解释性。

本文的主要工作和贡献包括：

\begin{enumerate}
    \item 提出了基于 Attention-MAPPO 的双层 MARL 框架，将轨迹控制与请求级卸载解耦为上层和下层两个协同决策层，降低动作空间复杂度。
    \item 设计了质量感知动作掩码和拉格朗日约束机制，分别处理不可行动作屏蔽和 DSR-MBS 负载权衡问题。
    \item 显式建模了 UAV 飞行能耗、UE 电池动态和 WPT 机制，维持系统长期运行的可持续性。
\end{enumerate}

实验结果表明：
\begin{enumerate}
    \item 完整双层 MARL 相比基础基线在 reward（+4482）、能耗（-56.6M）、公平性（+0.162）上取得统计显著改善，并将 offline rate 从 0.58\% 降至 0.00\%。
    \item 上层 attention-MAPPO 轨迹控制相对 uncoordinated greedy 对公平性和整体性能有独立且显著的贡献（$\Delta$reward +3805）；补充的 vanilla MAPPO 消融进一步表明，attention 机制主要体现为降低 MBS load 和能耗倾向，并改变轨迹协同与负载分配偏好。
    \item 能耗优先配置将 DSR、能耗和公平性之间的权衡显式化；通过 DSR-priority 配置可将 DSR 从 0.2083 提升至 0.2435，并在保持约 28\% 能耗节省的同时将其与 baseline 的差距缩小约 54\%，验证了框架的可调节性。
    \item Lagrange 约束有效控制 MBS load ratio（25.0\% $\rightarrow$ 13.7\%）。
    \item 质量感知 action mask 和 request attention 机制各自独立地改变了卸载分布和能耗/负载权衡模式。
    \item 端到端 \texttt{joint\_mappo} baseline 在相同训练预算下评估 reward 为 -9865.1、DSR 为 0.1349、fairness 为 0.4581，整体低于分层策略，说明直接联合轨迹-卸载动作学习存在样本效率和稳定性挑战。
\end{enumerate}

\section{未来展望}
\label{sec:future-work}

本文工作仍可从以下方向进一步扩展：

\begin{enumerate}
    \item 在更多 DSR 阈值和权重组合上进行细粒度 Pareto 扫描，以更好地理解 DSR-能耗-公平性之间的权衡关系。
    \item 扩大 vanilla/attention/joint MAPPO 补充实验的 training seeds 和 workload seeds，以获得与主实验一致的统计规模，提高实验结论的可靠性。
    \item 引入三维轨迹控制以更好地模拟真实部署环境，考虑 UAV 飞行高度的变化对通信和能耗的影响。
    \item 研究更高效的 Lagrange 乘子自适应更新策略，以提高约束优化的收敛速度和稳定性。
    \item 将服务缓存和内容分发纳入协同优化框架，实现轨迹、卸载、缓存和资源分配的全面协同优化。
    \item 探索异构 UAV 场景，考虑不同类型 UAV 的计算能力、通信能力和能耗特性差异。
    \item 研究动态拓扑下的多 UAV 协同问题，考虑 UAV 的加入和退出对系统性能的影响。
\end{enumerate}

\section{本章小结}
\label{sec:summary-chap06}

本章总结了本文的研究工作和主要贡献，并展望了未来的研究方向。首先回顾了本文提出的基于 Attention-MAPPO 的双层 MARL 框架及其主要创新点；然后总结了实验结果，验证了该框架在奖励、能耗、公平性和离线率等指标上的显著改善；最后提出了未来的研究方向，包括细粒度 Pareto 扫描、三维轨迹控制、Lagrange 乘子自适应更新、服务缓存和内容分发协同优化等。
